% Ported verbatim (prose unchanged) from docs/SECTION_2_DRAFT.md.
% All \cite{} keys resolved against the rebuilt paper/references.bib
% (26 entries, from docs/CITATIONS.md's human-verified 31-heading update).
% Two of the human's explicit sentence-level citation placements
% (missbench2026, farhadipour2025multimodal) are joined in the same
% sentence by additional newly-added entries with an equally strong or
% stronger topical fit (lucic2018gans, melis2017evaluation,
% pineau2021reproducibility) rather than left unused -- see the chat
% report for the reasoning.
\section{Related Work}
\label{sec:related-work}

\subsection{Graph and Contextual Architectures for Emotion Recognition in Conversation}

The architectural history of ERC is a history of adding conversational structure on top of utterance representations. Early multimodal systems classified utterances in isolation from concatenated or tensor-fused unimodal features~\cite{zadeh2017tfn}. Recurrent contextual models (DialogueRNN~\cite{dialoguernn2019} chief among them) introduced per-speaker state tracked across the dialogue; graph formulations generalized the idea, treating utterances or speakers as nodes and modeling their influence with graph convolutions and attention: DialogueGCN~\cite{dialoguegcn2019}, MMGCN~\cite{mmgcn2021} for the multimodal case, COGMEN~\cite{cogmen2022} with relational graph transformers, and a continuing line of refinements including shift-aware multi-task designs such as EmoShiftNet~\cite{emoshiftnet2025} and our own SocialArcNet~\cite{socialarcnet2026}, which maintains a recurrent per-speaker graph memory. Two properties of this literature matter for the present study. First, its experimental contract: with few exceptions, the contextual architecture is trained over frozen or lightly adapted encoder features: a pragmatic necessity in the era of expensive fine-tuning, and the regime our Section~\ref{sec:meld-frozen} reproduces, gains included. Second, its evaluation practice: components are typically assessed by scratch-retrained ablation, a design whose confounds our Section~\ref{sec:meld-mechanism} measures directly. We emphasize that nothing in our results impugns the original findings; on the foundations of their day, the gains were real. Our question is whether the contract still holds when the foundation is brought to current strength.

\subsection{Contextual Text Encoding: The Foundation Moves}

A parallel line asks how much of ERC is solved by the text encoder alone when it is fine-tuned and shown the dialogue. EmoBERTa~\cite{emoberta2021} demonstrated that a RoBERTa encoder given speaker names and preceding utterances in its input window reaches or exceeds contemporaneous full multimodal systems on MELD; CoMPM~\cite{compm2021} and successors~\cite{spcl2022} refined the recipe with speaker-aware memory and contrastive objectives. This line's implication for architecture research has, in our reading, been under-absorbed: it changes the baseline that any structural addition must beat, from a frozen-feature classifier to a task-adapted contextual encoder several points stronger. Our contribution is to measure that implication directly and under controls (the same components, attributed against both foundations with one instrument) rather than to infer it from cross-paper comparisons, which differ in splits, features, and tuning budgets.

\subsection{Parameter-Efficient Adaptation}

Low-Rank Adaptation~\cite{hu2022lora} and the broader parameter-efficient fine-tuning family make the strong-foundation regime accessible at commodity cost: our text foundations train on a single consumer GPU in minutes per epoch, which is precisely why the frozen-feature contract can no longer be justified by resource constraints alone in most settings. PEFT has been applied within affective computing for domain and speaker adaptation~\cite{feng2023peftser}; we use it not as a contribution but as the enabling condition of the study: the boundary we chart exists because adaptation became cheap. Our results also carry one practical PEFT-adjacent finding (Section~\ref{sec:meld-foundation}): training-time logit adjustment interacts badly with this recipe in a way its inference-time form does not.

\subsection{Ablation Validity and the Reporting of Negative Results}

Our methodological stance draws on a growing literature questioning what ablation studies measure: critiques of unequal optimization budgets between compared configurations~\cite{missbench2026,lucic2018gans,melis2017evaluation}, demonstrations that simple baselines match elaborate architectures when tuned comparably~\cite{liang2019strongsimple}, and calls for pre-registration and negative-result reporting in machine learning~\cite{farhadipour2025multimodal,pineau2021reproducibility}. The residual attribution construction of Section~\ref{sec:residual-construction} is our instrument-level response: rather than equalizing tuning effort statistically, it removes the re-learning confound by construction, at the disclosed cost of a soft prior toward the foundation. To our knowledge its use as an attribution instrument for conversation-level architecture is new, though the ingredient (zero-initialized residual branches) has precedent in normalization-free networks and adapter methods~\cite{bachlechner2020rezero,zhang2019fixup}. In spirit, the closest relatives of this paper are studies asking whether structural machinery survives stronger foundations in adjacent fields (e.g., whether graph structure helps over strong feature encoders in node classification~\cite{zhang2021graphless,huang2020labelprop}), and we read our results as ERC's instance of that recurring pattern: structure compensates for representation until representation catches up.
